\documentclass{article}

\usepackage{booktabs}
\usepackage{tabularx}
\usepackage{hyperref}

\hypersetup{
    colorlinks=true,       % false: boxed links; true: colored links
    linkcolor=red,          % color of internal links (change box color with linkbordercolor)
    citecolor=green,        % color of links to bibliography
    filecolor=magenta,      % color of file links
    urlcolor=cyan           % color of external links
}

\title{Hazard Analysis\\\progname}

\author{\authname}

\date{}

\input{../Comments.text}
\input{../Common.text}

\begin{document}

\maketitle
\thispagestyle{empty}

~\newpage

\pagenumbering{roman}

\begin{table}[hp]
\caption{Revision History} \label{TblRevisionHistory}
\begin{tabularx}{\textwidth}{llX}
\toprule
\textbf{Date} & \textbf{Developer(s)} & \textbf{Change}\\
\midrule
Date1 & Name(s) & Description of changes\\
Date2 & Name(s) & Description of changes\\
... & ... & ...\\
\bottomrule
\end{tabularx}
\end{table}

~\newpage

\tableofcontents

~\newpage

\pagenumbering{arabic}

\wss{You are free to modify this template.}

\section{Introduction}

\wss{You can include your definition of what a hazard is here.}

\section{Scope and Purpose of Hazard Analysis}

\wss{You should say what \textbf{loss} could be incurred because of the
hazards.}

\section{System Boundaries and Components}

\wss{Dividing the system into components will help you brainstorm the hazards.
You shouldn't do a full design of the components, just get a feel for the major
ones.  For projects that involve hardware, the components will typically include
each individual piece of hardware.  If your software will have a database, or an
important library, these are also potential components.}

\section{Critical Assumptions}

\wss{These assumptions that are made about the software or system.  You should
minimize the number of assumptions that remove potential hazards.  For instance,
you could assume a part will never fail, but it is generally better to include
this potential failure mode.}

\section{Failure Mode and Effect Analysis}

\wss{Include your FMEA table here. This is the most important part of this document.}
\wss{The safety requirements in the table do not have to have the prefix SR.
The most important thing is to show traceability to your SRS. You might trace to
requirements you have already written, or you might need to add new
requirements.}
\wss{If no safety requirement can be devised, other mitigation strategies can be
entered in the table, including strategies involving providing additional
documentation, and/or test cases.}

\section{Safety and Security Requirements}

The following security, safety, and reliability requirements were drawn as a result of the hazard analysis process:

\begin{itemize}
\item \textbf{FR-HA1}: The system shall prompt the user when noisy (less than 15 dB) microphone inputs are passed in, and allow them to retry that input.\\
\textbf{Rationale:} The system should not take in corrupted data as it will degrade the quality of the results and misrepresent a patient’s standing. A safer alternative is to prompt the user of the error, reject the data, and offer a retry option when the issue is resolved. 15 dB is the standard signal to noise ratio aimed for a clear microphone signal \textcolor{red}{[ACCEPTABLE MIC NOISE]}.

\item \textbf{FR-HA2}: The system must back up all patient progress data and analytics at the end of each session. It must alert the user if the backup fails.\\
\textbf{Rationale:} In the event of a data breach, patient data is not lost.

\item \textbf{FR-HA3}: The speech engine must provide a confidence percentage for each phoneme deduction it reaches. (have a FR to connect this to)\\
\textbf{Rationale:} Upon some speech analysis, the results must be kept transparent so the SLP is confident to accept them. This way they can leverage their professional judgement to reject incorrect deductions and consider high confidence ones.

\item \textbf{FR-HA4}: The system must alert the user when WiFi connection becomes unstable. It must provide an offline version which records patient audio for when the WiFi recovers.\\
\textbf{Rationale:} If the network connection becomes unstable, the connection to any processing server may drop and the analysis system may become inactive. The SLP needs to know of this so they could take alternative measures instead of randomly finding out their session was not tracked. A metric for unstable will depend on the user’s hardware specifications and the server specifications.

\item \textbf{FR-HA5}: The system must be able to transfer from online to offline mode at all times. The transfer cannot be blocked by network or interface malfunctions.\\
\textbf{Rationale:} If the network connection drops, the system cannot freeze and must fail gracefully. This way, the session may continue with limited functionality while still allowing the SLP to save speech files, and access the interface.

\item \textbf{FR-HA6}: The system must have two factor authentication and a password recovery system in place.

\item \textbf{FR-HA7}: The system must preface all deductions and analytics with an acknowledgement that they may be incorrect and may only be accepted under the discretion of an SLP.\\
\textbf{Rationale:} The tool is purely assistive and software may not be 100\% correct. Therefore this disclosure is necessary to ensure users do not try to replace the judgement of SLPs with this.

\item \textbf{FR-HA8}: Prior to authorizing any audio recordings, a one-time confirmation code sent to the patient’s trusted device at the start of the session must be entered into the SLP interface.\\
\textbf{Rationale:} This serves as a method of consent for the clinic to record the patient during this session. It indirectly verifies that the SLP has their interface open to the correct patient as well.

\item \textbf{FR-HA9}: The system shall maintain an audit trail detailing the user, action, and time of all data accesses and manipulations.\\
\textbf{Rationale:} It is important to keep a traceable log of all user and system actions in the event of errors, data breaches or other undesired events.
\end{itemize}

\section{Roadmap}

Requirements 1, 3, 7 will be integrated within the capstone timeline. The core of the product is its speech analysis ability, so filtering inputs for noisy data is crucial. It is important to be transparent about the validity of a claim, which is why each deduction must be supported by a confidence level. This way an SLP can make a more informed judgement rather than blindly trusting the software. Full disclosure is needed at all times to ensure it does not impede the professional judgement of an SLP. 

The remaining requirements specify necessary security and reliability mechanisms that must be in place before the product is released. Elements like authentication and redundancy are especially crucial for real life clinic scenarios, but do not enhance the value of the project as a capstone. Instead, time is better spent enhancing the SLP interface of the system, and focusing on the speech analysis. 

\newpage{}

\section{Appendix --- Generative AI Use Disclosure}

\wss{ChatGPT (version 5) was used to brainstorm ideas for the template for this disclose section.}

\subsection{AI Tools Used}

\wss{Give the name of the AI tool(s) used and the version.}

\subsection{How and Where Used}

\wss{Explain what the tool(s) were used for.  Examples include: brainstorming,
explaining a technical concept, reviewing the document, generating or modifying
text, generating or modifying diagrams, identifying errors or omissions or
inconsistencies.}

\wss{For each of the uses, identify which parts of the document were affected.}

\wss{You do not need to report every interaction, but you should disclose the 
uses that materially contributed to your work}

\subsection{Verification of AI-Generated Content}

\wss{\begin{itemize}
    \item Describe how the team checked AI-generated or AI-assisted material
    before including it in the document.
    \item In particular, verify factual claims, technical assertions,
    requirements, calculations, references, and other information for which an
    AI system may produce incorrect or fabricated results.
    \item \textbf{The team is responsible for the accuracy, completeness,
consistency, and appropriateness of the submitted document, regardless of
whether material was generated by a human or an AI system.}
\end{itemize}}

\newpage{}

\section{Appendix --- Reflection}

\wss{Not required for CAS 741}

\input{../Reflection.text}

\begin{enumerate}
    \item What went well while writing this deliverable? 
    \item What pain points did you experience during this deliverable, and how
    did you resolve them?
    \item Which of your listed risks had your team thought of before this
    deliverable, and which did you think of while doing this deliverable? For
    the latter ones (ones you thought of while doing the Hazard Analysis), how
    did they come about?
    \item Other than the risk of physical harm (some projects may not have any
    appreciable risks of this form), list at least 2 other types of risk in
    software products. Why are they important to consider?
\end{enumerate}

\end{document}